% ============ 第1章 引言 ============
\chapter{引言}

\section{背景}

Wi-Fi 信道状态信息（Channel State Information, CSI）刻画了无线信号从发送端到接收端经历的多径叠加结果：环境中人或物体的移动会改变多径传播，进而引起 CSI 幅值与相位的波动。利用这种“无需被测设备携带任何传感器”的无线感知（Wi-Fi Sensing）能力，可以用普通的商用 Wi-Fi 芯片实现运动检测、存在检测等应用\cite{ref:hernandez}。IEEE 于 2025 年 5 月正式批准了 802.11bf WLAN Sensing 标准\cite{ref:80211bf}，标志着该方向已从学术研究进入行业标准化阶段。

在工程侧，乐鑫（Espressif）官方开源工程 esp-csi 为 ESP32 系列芯片提供了 CSI 采集与检测组件（esp-radar）\cite{ref:espcsi}，使廉价微控制器可以承担边缘感知任务\cite{ref:hernandez}。社区实践（如 ESPectre 运动检测框架\cite{ref:espectre}）与学术研究同时指出两个关键约束：其一，ESP32 的 CSI 不携带 AGC 增益信息，射频前端增益切换会造成乘性幅值突变，需在检测端加以过滤\cite{ref:agc185,ref:csisim}；其二，商品级 ESP32 在多人/多目标分辨上是公认弱项，多节点系统应以“每条链路独立检测”管理预期\cite{ref:gait}。

\section{选题与任务}

本组选题为课程大作业题目3：\textbf{基于 ESP32 CSI 扰动特征的多监测点监控系统}。题目要求概述如下：发送节点分布于各监测区，一块 ESP32-S3 接收节点利用 esp-radar 组件的多 peer 能力按源 MAC 分发 CSI，对每条链路独立完成运动/存在检测，并经 UART 二进制协议上报；PC 端 Python 桌面监控程序展示区域状态卡片、实时特征曲线与事件日志，支持 CSV 记录与回放。

本组对该题目的理解为：这是一个“多节点无线感知 + 嵌入式实时检测 + 串行通信 + 上位机可视化”的综合系统设计题，核心难点在于（1）多发送节点共享同一信道时，接收端必须能按来源区分链路并互不串扰；（2）检测算法须在廉价微控制器上实时运行，且对环境变化可现场标定；（3）上行链路须在控制台日志与二进制帧交织的真实串口流上保证数据完整性；（4）作为微机原理课程大作业，系统须覆盖串行通信、中断与回调、定时器、存储器组织、多任务与并发保护等课程知识点。

\section{本文工作与组织}

本文完成了题目3系统的设计、实现与台架联调。主要工作包括：搭建“1 收 3 发 + 路由器组网 + 直收兜底”的系统拓扑；设计带 CRC 校验与失步重同步的 UART 二进制上报协议；实现逐链路滑动窗统计与滞回状态机检测算法（含 RSSI 突变帧过滤）；实现现场标定与 NVS 持久化；开发 PySide6 桌面上位机；完成协议自测、数据流、帧质量、双链路联调与上位机冒烟共五组台架实验。

本文组织如下：第2章给出设计指标；第3章介绍硬件平台与开发环境；第4章阐述系统设计（含微机原理知识点映射）；第5章介绍关键模块实现；第6章汇总已完成的台架测试与验证；第7章为现场部署与验收安排（待执行）；第8章总结全文并讨论局限；第9章为小组分工；附录A/B 给出协议帧格式总表与工程文件索引。
